% !TeX root = ../papers/nc-p3.tex
\section{非正則な水平曲線の種数3水準への特殊化}
\label{sec:rank-fifteen-horizontal-specialization}
\begin{theorem}\label{thm:rank-fifteen-horizontal-specialization}
各階数15表示について、命題\ref{prop:rank-fifteen-pencil-plane}の三二次式を
$c$ の相対射影曲線とし、窓の表示のため $L^2=1+c^2$ を選ぶ。
この曲面の既約曲線 $B$ が $c$ の直線を支配し、その幾何学的一般窓が通常かつ
標準Koszul左右AS正則でないと仮定する。$c_0^2-c_0+1=0$ の各水準で
$B$ の射影閉包が取り得る特殊化は、$K=\pm1$ 上の八つの滑らかな点に限られる。特に各 $c_0$ 上で閉包はこの八点のいずれかを通る。
\end{theorem}
\begin{proof}
$c_0$ で $s=1$、$1+c_0^2\ne0$ である。$B$ の閉包の正規化上で
$c_0$ の上に付値を取る。曲線族は $c_0$ の近傍で射影的であるから、
この特殊化は存在する。必要なら有限拡大して幾何点と $L$ の選択を取る。

特殊点で $K$ が通常なトレースなら、三二次式の曲線は滑らかで
八次モデルの通常窓を与える。定理\ref{thm:rank-fifteen-octic-center}により
経路長4,5と四次中心の次元は周囲の通常部分の一般次元に一致する。
定理\ref{thm:rank-fifteen-octic-all}の証明はこの特殊点に中心四式の有限商の
証明書を与える。道・中心の定階数性により、付値環の一般窓にこの証明書を持ち上げられる。
従って $B$ の仮定に反する。

通常でない場合も $K=\infty$ なら定理\ref{thm:rank-fifteen-pencil-infinity}の
重みを使う。先頭係数は $c,L$ の有理式で、同定理の検算により $c_0$ の上で単元である。
従って $c$ も変動する付値環で同じ重みと剰余関係が使える。
$K^2+K-1=0$ の各有限点では、命題\ref{prop:rank-fifteen-octic-residues}が
実際の関係を評価した $r,w,z-1,w^2-1$ が単元のチャートが使える。
そこでも剰余の道と中心の次元および四中心商の証明書が得られており、
これらの単元は曲面上の特殊化でも保たれる。

$K=\pm2$ の八点は定理\ref{thm:rank-fifteen-pencil-ramified-ends}の
$\delta=\epsilon$ の二分岐に当たる。$q$ の二根は単純で
$E'(K_0)=-3$、$n_0\ne0$ なので、同じ相対座標 $u=x-\epsilon$ が使える。
\path{computations/verify_rank_fifteen_horizontal_specialization.py} は109における
$c^2-c+1$ の二根、二つの分岐、両表示および各四シートで関係の各文字ブロックを調べる。
定数項が全て零なら、固定した水準での零に加えて
$\mathbb Q(c)[q]/(q^2-2(2K_0-\epsilon)q+1)$ における恒等的消滅を確認する。
従って曲面でも $u$ で割ることができる。一度以内の除算後に単元ピボットがあり、
得られた剰余上で道・中心が所要の次元を持ち、全四始点で商の長さ13までの上界を満たす。
このため $c_0$ 上のこれらの点を中心とする付値でも一般窓が正則になる。

さらに $K=0$ では $H=\pm i$、$A=B=0$ である。特殊ファイバーの
$A^2=K(K+2)$、$B^2=K(K-2)$ の接錐は相異なる二直線で、ここには二つの節点がある。
一方、\emph{曲面上}では $u=r+1,v=w-1$ が正則局所座標になる。実際
\[
 R=(r+r^{-1})/2,\quad W=(w+w^{-1})/2,\quad
 K=(R+W)/2,\quad s=2(K^2-1)/(R-W).
\]
中心で $R=-1,W=1$ なので分母は単元である。
$c+c^{-1}=s$ と $n^2=(K^2-1)/(1+c^2)$ はいずれも境界値において単根で、
$n_0=\pm c_0$ となる。$H=Ln$ も単元であり、$A=s(r-R)/H$、$B=s(w-W)/H$ で戻せる。
$R,W$ はともに一次微分が零なので、$K,c,n$ の二方向の一次微分も零である。

上記の検算は二つの節点と両表示について各Pl\"ucker座標の組を評価する。
単元がない組に対し、$u=0$ の分岐 $(\delta,\epsilon)=(1,-1)$、
または $v=0$ の分岐 $(-1,1)$ で全座標が恒等的に零となることを、先の二次式を法として
確かめる。この二つの因子は正則座標そのものなので、対応する $u$ または $v$ で
割れる。先頭値はそれぞれ $r$ または $w$ での偏微分である。両定数根 $c_0$ と両表示で、
各節点につき二組を $u$、別の二組を $v$ で割ると全組が単元座標を持つ。
得られる関係で道・中心・四中心商の証明書を再び確認した。
この関係格子は独立な二変数を持つので、節点を通る任意の弧の通常一般窓に使える。

最後に $K=\pm1$ の上には $H=0$ で他の二根が非零となる点が各四つあり、いずれも滑らかである。
以上で他の全ての特殊化を除いた。
\end{proof}
